%% toc: P2 21 21.2,21.3
\section{Proofs for Craig interpolation with free variables and Beth definability}
\label{sec:annex-p02-w02-p02-037-craig-beth}

Throughout, signatures, formulas, substitution, alpha-equivalence and
satisfaction are those of \cref{sec:fol-syntax-semantics-01},
\cref{sec:fol-syntax-semantics-02}, \cref{sec:fol-syntax-semantics-04} and
\cref{sec:fol-semantics-02}.  Structures have nonempty domains.  For a
formula \(\eta\), \(\operatorname{Var}(\eta)\) is the set of all variables
occurring in \(\eta\), bound or free, and \(\operatorname{voc}(\eta)\) is its
occurrence vocabulary (Convention~\ref{def:p02-w02-p02-037-vocabulary}).

\begin{annexitem}{annex:p02-w02-p02-037-craig-formula}{thm:p02-w02-p02-037-craig-formula}{Craig interpolation, formula form}
\begin{proof}
Put \(\Lang=\Lang_1\cup\Lang_2\).  Let
\(V=\FV(\varphi)\cup\FV(\psi)=\{v_1,\ldots,v_k\}\) with the \(v_i\)
pairwise distinct, choose pairwise distinct constant symbols
\(c_1,\ldots,c_k\notin\Lang\), and put
\(\Lang^{\mathrm{c}}=\Lang\cup\{c_1,\ldots,c_k\}\).  If \(k=0\) the constant list is
empty and steps~A to~C below are vacuous.

\paragraph{A. Substituting a constant.}
Let \(\eta\) be a formula, \(x\) a variable and \(c\) a constant symbol.
Since \(\operatorname{Var}(c)=\varnothing\), the third quantifier clause of
capture-avoiding substitution never applies in computing \(\eta[c/x]\).
Induction along the clauses of that definition therefore shows that
\(\eta[c/x]\) is \(\eta\) with each free occurrence of \(x\) replaced by
\(c\) and with every quantifier unchanged, and that
\[
 \FV(\eta[c/x])=\FV(\eta)\setminus\{x\},
 \qquad
 \operatorname{voc}(\eta[c/x])=
 \begin{cases}
  \operatorname{voc}(\eta)\cup\{c\}, & x\in\FV(\eta),\\
  \operatorname{voc}(\eta), & x\notin\FV(\eta).
 \end{cases}
\]
In the quantifier case \(Qy\,\eta_0\) with \(y=x\), the formula is unchanged
and \(x\notin\FV(Qx\,\eta_0)\); with \(y\neq x\), the free variables are
\((\FV(\eta_0)\setminus\{x\})\setminus\{y\}\) and \(x\) is free in
\(Qy\,\eta_0\) exactly when it is free in \(\eta_0\).

Define
\(\varphi^{\mathrm{c}}=\varphi[c_1/v_1]\cdots[c_k/v_k]\) and
\(\psi^{\mathrm{c}}=\psi[c_1/v_1]\cdots[c_k/v_k]\).  Applying the displayed facts
\(k\) times, and using that the \(c_i\) are distinct and lie outside
\(\Lang\),
\[
\begin{gathered}
 \FV(\varphi^{\mathrm{c}})=\FV(\psi^{\mathrm{c}})=\varnothing,\\
 \operatorname{voc}(\varphi^{\mathrm{c}})=
   \operatorname{voc}(\varphi)\cup\{c_i:v_i\in\FV(\varphi)\},
\end{gathered}
\]
and likewise for \(\psi^{\mathrm{c}}\).  Hence \(\varphi^{\mathrm{c}}\) and
\(\psi^{\mathrm{c}}\) are \(\Lang^{\mathrm{c}}\)-sentences and
\[
\begin{aligned}
 &\operatorname{voc}(\varphi^{\mathrm{c}})\cap\operatorname{voc}(\psi^{\mathrm{c}})\\
 &\quad=\bigl(\operatorname{voc}(\varphi)\cap\operatorname{voc}(\psi)\bigr)
  \cup\{c_i:v_i\in\FV(\varphi)\cap\FV(\psi)\}.
\end{aligned}
 \tag{A}
\]

\paragraph{B. Satisfaction of a constant instance.}
For an \(\Lang^{\mathrm{c}}\)-structure \(\mathcal N\) and an assignment \(s\), put
\(s^{\mathcal N}=s[v_1\mapsto c_1^{\mathcal N}]\cdots
[v_k\mapsto c_k^{\mathcal N}]\).  Because the \(v_i\) are distinct, these
updates commute, so \(s^{\mathcal N}\) does not depend on their order.  For
every \(\Lang^{\mathrm{c}}\)-formula \(\eta\) and every listing
\(i_1,\ldots,i_k\) of \(1,\ldots,k\),
\[
 \mathcal N,s\entails\eta[c_{i_1}/v_{i_1}]\cdots[c_{i_k}/v_{i_k}]
 \quad\Longleftrightarrow\quad
 \mathcal N,s^{\mathcal N}\entails\eta .
 \tag{B}
\]
Indeed, the value of the closed term \(c_i\) is \(c_i^{\mathcal N}\) under
every assignment, so the substitution lemma, \cref{thm:fol-substitution},
applied to the outermost substitution gives
\(\mathcal N,s\entails\eta'[c_{i_k}/v_{i_k}]\) iff
\(\mathcal N,s[v_{i_k}\mapsto c_{i_k}^{\mathcal N}]\entails\eta'\); repeating
this \(k\) times with the updated assignments yields (B).

\paragraph{C. The sentence implication is valid.}
Let \(\mathcal N\) be an \(\Lang^{\mathrm{c}}\)-structure with
\(\mathcal N\entails\varphi^{\mathrm{c}}\), and let \(s\) be any assignment.  By (B),
\(\mathcal N,s^{\mathcal N}\entails\varphi\).  By
\cref{lem:reduct-invariance},
\(\mathcal N\mathbin{\upharpoonright}\Lang,s^{\mathcal N}\entails\varphi\).  The hypothesis
\(\entails\varphi\to\psi\) over \(\Lang\)-structures gives
\(\mathcal N\mathbin{\upharpoonright}\Lang,s^{\mathcal N}\entails\psi\), hence
\(\mathcal N,s^{\mathcal N}\entails\psi\) by
\cref{lem:reduct-invariance} again, and
\(\mathcal N\entails\psi^{\mathrm{c}}\) by (B).  Thus
\(\entails\varphi^{\mathrm{c}}\to\psi^{\mathrm{c}}\).

\paragraph{D. A sentence interpolant.}
The sentences \(\varphi^{\mathrm{c}}\) and \(\psi^{\mathrm{c}}\) lie in the compatible
signatures \(\Lang_1\cup\{c_1,\ldots,c_k\}\) and
\(\Lang_2\cup\{c_1,\ldots,c_k\}\).  The sentence case of
Theorem~\ref{thm:craig-interpolation} applies to them.  Its proof in
Annex~\ref{annex:p7.interp.3} reaches sentence interpolation in a general
signature before, and independently of, its paragraph on free variables.
The same sentence case is also proved by the route through joint
consistency in Annex~\ref{annex:p3.abstract.5}, using
Theorem~\ref{thm:p03-robinson-joint-consistency}.  Either proof gives an
\(\Lang^{\mathrm{c}}\)-sentence \(\theta\) with
\[
\begin{gathered}
 \entails\varphi^{\mathrm{c}}\to\theta,
 \qquad
 \entails\theta\to\psi^{\mathrm{c}},\\
 \operatorname{voc}(\theta)\subseteq
   \operatorname{voc}(\varphi^{\mathrm{c}})\cap\operatorname{voc}(\psi^{\mathrm{c}}),
\end{gathered}
\]
where equality and \(\bottom\) may occur in \(\theta\).  By (A), a
constant \(c_i\) occurs in \(\theta\) only if
\(v_i\in\FV(\varphi)\cap\FV(\psi)\), and every other symbol of
\(\operatorname{voc}(\theta)\) lies in
\(\operatorname{voc}(\varphi)\cap\operatorname{voc}(\psi)\).

\paragraph{E. Renaming bound variables away from \(V\).}
\emph{Claim.}  For every formula \(\eta\) and finite set \(W\) of variables
there is a formula \(\eta^{\circ}\) with
\(\eta^{\circ}\equiv_\alpha\eta\),
\(\FV(\eta^{\circ})=\FV(\eta)\),
\(\operatorname{voc}(\eta^{\circ})=\operatorname{voc}(\eta)\), and no
quantifier of \(\eta^{\circ}\) binding a member of \(W\).

Induct on \(\eta\).  Atomic formulas have no quantifier; take
\(\eta^{\circ}=\eta\).  For \(\bottom\) and the connectives, apply the
induction hypothesis to the immediate subformulas and use that
\(\equiv_\alpha\) is a congruence.  Let \(\eta=Qy\,\eta_0\), and let
\(\eta_0^{\circ}\) be given by the induction hypothesis.  If
\(y\notin W\), put \(\eta^{\circ}=Qy\,\eta_0^{\circ}\).  If \(y\in W\),
choose a variable \(z\notin\operatorname{Var}(\eta_0^{\circ})\cup W\cup\{y\}\)
and put \(\eta^{\circ}=Qz\,(\eta_0^{\circ}[z/y])\).  Then
\(Qy\,\eta_0\equiv_\alpha Qy\,\eta_0^{\circ}\) by congruence and
\(Qy\,\eta_0^{\circ}\equiv_\alpha Qz\,(\eta_0^{\circ}[z/y])\) is a generating
step, since \(z\notin\operatorname{Var}(\eta_0^{\circ})\).  In computing
\(\eta_0^{\circ}[z/y]\) the renaming clause never applies, because no
quantifier of \(\eta_0^{\circ}\) binds \(z\), and no quantifier of
\(\eta_0^{\circ}\) binds \(y\in W\); so \(\eta_0^{\circ}[z/y]\) is
\(\eta_0^{\circ}\) with every occurrence of \(y\) replaced by \(z\), its
quantifiers are those of \(\eta_0^{\circ}\), and its vocabulary is
unchanged.  As \(z\notin\FV(\eta_0^{\circ})\),
\[
 \FV(Qz\,(\eta_0^{\circ}[z/y]))
 =\bigl((\FV(\eta_0^{\circ})\setminus\{y\})\cup
   \{z: y\in\FV(\eta_0^{\circ})\}\bigr)\setminus\{z\}
 =\FV(\eta_0)\setminus\{y\}.
\]
This proves the claim.

\emph{Claim.}  If no quantifier of \(\eta\) binds \(v\) and
\(v\notin\FV(\eta)\), then \(v\notin\operatorname{Var}(\eta)\).  For atomic
\(\eta\), every variable occurrence is free.  The connective cases follow
from the induction hypothesis.  For \(Qy\,\eta_0\) we have \(y\neq v\), so
\(v\notin\FV(\eta_0)\), and the induction hypothesis gives
\(v\notin\operatorname{Var}(\eta_0)\); hence
\(v\notin\operatorname{Var}(\eta_0)\cup\{y\}\).

Apply the first claim to \(\theta\) with \(W=V\), obtaining
\(\theta^{\circ}\).  It is a sentence with the vocabulary of \(\theta\), and
by the second claim no \(v_i\) occurs in it.  By alpha-invariance,
\cref{cor:alpha-invariance}, \(\theta^{\circ}\) and \(\theta\) are satisfied
by the same structures and assignments.

\paragraph{F. Replacing constants by variables.}
For a term \(t\), a constant \(c\) and a variable \(v\), let
\(t\{v/c\}\) replace every occurrence of \(c\) in \(t\) by \(v\).  For a
formula \(\eta\), let \(\eta\{v/c\}\) apply this replacement to every term
of every atomic subformula and leave \(\bottom\), the connectives and every
quantifier with its variable unchanged.  If
\(v\notin\operatorname{Var}(\eta)\), then
\[
\begin{gathered}
 \eta\{v/c\}[c/v]=\eta,
 \qquad
 \operatorname{voc}(\eta\{v/c\})=\operatorname{voc}(\eta)\setminus\{c\},\\
 \FV(\eta\{v/c\})=\FV(\eta)\cup\{v:c\in\operatorname{voc}(\eta)\}.
\end{gathered}
\]
For terms, induction gives \(t\{v/c\}[c/v]=t\) when \(v\notin\operatorname{Var}(t)\):
the constant \(c\) goes to \(v\) and back, a variable \(y\neq v\) is fixed
by both operations, and compound terms are handled argumentwise.  For
formulas, the atomic and connective cases follow from the term case and the
induction hypothesis.  For \(Qy\,\eta_0\), the hypothesis gives \(y\neq v\),
and \(y\notin\operatorname{Var}(c)=\varnothing\), so the second quantifier
clause of substitution applies:
\((Qy\,\eta_0\{v/c\})[c/v]=Qy\,(\eta_0\{v/c\}[c/v])=Qy\,\eta_0\).  The
free-variable identity holds because each new occurrence of \(v\) stands
where \(c\) stood and no quantifier binds \(v\); in the quantifier case
\(y\neq v\) keeps \(v\) outside the removed set \(\{y\}\).  The vocabulary
identity is immediate, since only occurrences of \(c\) change.

Now put \(\eta_0=\theta^{\circ}\) and
\(\eta_i=\eta_{i-1}\{v_i/c_i\}\) for \(i=1,\ldots,k\), and let
\(\chi=\eta_k\).  At stage \(i\),
\(\operatorname{Var}(\eta_{i-1})\subseteq
\operatorname{Var}(\theta^{\circ})\cup\{v_1,\ldots,v_{i-1}\}\), which does not
contain \(v_i\).  So the displayed identities apply at every stage, and
\[
 \chi[c_k/v_k][c_{k-1}/v_{k-1}]\cdots[c_1/v_1]=\theta^{\circ},
\]
\[
\begin{aligned}
 \operatorname{voc}(\chi)
   &=\operatorname{voc}(\theta)\setminus\{c_1,\ldots,c_k\}
   \subseteq\operatorname{voc}(\varphi)\cap\operatorname{voc}(\psi),\\
 \FV(\chi)&=\{v_i:c_i\in\operatorname{voc}(\theta)\}
   \subseteq\FV(\varphi)\cap\FV(\psi),
\end{aligned}
\]
using step~D for both inclusions and \(\FV(\theta^{\circ})=\varnothing\).
In particular \(\chi\in\operatorname{Fm}(\Lang)\).

\paragraph{G. The two implications.}
Let \(\mathcal M\) be an \(\Lang\)-structure and \(s\) an assignment.  Let
\(\mathcal N\) be the \(\Lang^{\mathrm{c}}\)-expansion of \(\mathcal M\) with
\(c_i^{\mathcal N}=s(v_i)\) for every \(i\).  Then
\(s^{\mathcal N}=s\).  Different constants may receive the same value; no
step below uses distinctness of values.

Suppose \(\mathcal M,s\entails\varphi\).  By \cref{lem:reduct-invariance},
\(\mathcal N,s^{\mathcal N}\entails\varphi\), so
\(\mathcal N\entails\varphi^{\mathrm{c}}\) by (B), and
\(\mathcal N\entails\theta\) by step~D.  By alpha-invariance
\(\mathcal N,s\entails\theta^{\circ}\), that is,
\(\mathcal N,s\entails\chi[c_k/v_k]\cdots[c_1/v_1]\).  By (B),
\(\mathcal N,s^{\mathcal N}\entails\chi\), and since \(s^{\mathcal N}=s\) and
\(\chi\in\operatorname{Fm}(\Lang)\), \cref{lem:reduct-invariance} gives
\(\mathcal M,s\entails\chi\).  Hence \(\entails\varphi\to\chi\).

Suppose \(\mathcal M,s\entails\chi\).  The same steps in reverse give
\(\mathcal N,s\entails\theta^{\circ}\), so \(\mathcal N\entails\theta\) and
\(\mathcal N\entails\psi^{\mathrm{c}}\) by step~D; then (B) gives
\(\mathcal N,s^{\mathcal N}\entails\psi\), and \cref{lem:reduct-invariance}
gives \(\mathcal M,s\entails\psi\).  Hence \(\entails\chi\to\psi\).

Steps~F and~G give the vocabulary and free-variable conditions and both
implications.  When \(\varphi\) and \(\psi\) are sentences, \(k=0\) and
\(\chi=\theta^{\circ}\) is a sentence.
\end{proof}
\end{annexitem}

\begin{annexitem}{annex:p02-w02-p02-037-equality-needed}{thm:p02-w02-p02-037-equality-needed}{Equality is needed in interpolants}
\begin{proof}
If \(\mathcal M,s\entails\varphi\), there are \(a,b\) with
\(a\in P^{\mathcal M}\) and \(b\notin P^{\mathcal M}\); then \(a\neq b\), and
the assignment sending \(x,y\) to \(a,b\) satisfies \(\neg x=y\).  Hence
\(\entails\varphi\to\psi\).  Clearly \(\operatorname{voc}(\varphi)=\{P\}\)
and \(\operatorname{voc}(\psi)=\varnothing\).

Let \(\chi\) contain no equality and have
\(\operatorname{voc}(\chi)=\varnothing\).  An atomic formula is an equation
or a relation atom; the first kind is excluded and the second would put its
relation symbol into \(\operatorname{voc}(\chi)\).  So the only atomic
subformula of \(\chi\) is \(\bottom\).  By induction on \(\chi\) there is a
truth value \(\tau(\chi)\) such that, for every structure \(\mathcal M\) and
every assignment \(s\), \(\mathcal M,s\entails\chi\) iff \(\tau(\chi)\) is
truth.  For \(\bottom\) take falsity.  For the connectives combine the
values by their truth clauses.  For \(\forall x\,\chi_0\) and
\(\exists x\,\chi_0\), every update \(s[x\mapsto a]\) satisfies \(\chi_0\)
iff \(\tau(\chi_0)\) is truth; since the domain is nonempty, at least one
update exists, so both quantified formulas have the value
\(\tau(\chi_0)\).

If \(\tau(\chi)\) is truth, then \(\entails\chi\to\psi\) would make
\(\psi\) valid.  But in a structure with one element every assignment gives
\(x\) and \(y\) the same value, so \(\psi\) is false there.  If
\(\tau(\chi)\) is falsity, then \(\entails\varphi\to\chi\) would make
\(\varphi\) unsatisfiable.  But the structure with domain \(\{0,1\}\) and
\(P=\{0\}\) satisfies \(\varphi\).  In both cases one of the two
implications fails.
\end{proof}
\end{annexitem}

\begin{annexitem}{annex:p02-w02-p02-037-renaming}{thm:p02-w02-p02-037-renaming}{Renaming the predicate}
\begin{proof}
The structures \(\mathcal N\) and \(\mathcal N^{\mathrm{ren}}\) have the same domain
and interpret every function symbol of \(\Lang\) identically.  The terms of
\(\Lang(P)\) are \(\Lang\)-terms, because \(P\) is a relation symbol, so
induction on terms gives
\(\sem{t}^{\mathcal N}_s=\sem{t}^{\mathcal N^{\mathrm{ren}}}_s\) for every term
\(t\) and assignment \(s\).

Induct on \(\alpha\).  For an equation \(t=u\) and for a relation atom with
a symbol of \(\Lang\), the formula \(\alpha'\) is \(\alpha\), the term values
agree, and the two structures interpret the symbol identically.  For
\(\alpha=P(t_1,\ldots,t_n)\), \(\alpha'=P'(t_1,\ldots,t_n)\), and
\[
\begin{aligned}
 \mathcal N,s\entails P'(t_1,\ldots,t_n)
 &\iff
 (\sem{t_1}^{\mathcal N}_s,\ldots,\sem{t_n}^{\mathcal N}_s)\in P'^{\mathcal N}
 =P^{\mathcal N^{\mathrm{ren}}}\\
 &\iff
 \mathcal N^{\mathrm{ren}},s\entails P(t_1,\ldots,t_n);
\end{aligned}
\]
for \(n=0\) the tuple is empty and the same equivalence holds.  Both
structures falsify \(\bottom\).  Renaming commutes with the connectives, so
those cases follow from the induction hypothesis.  Renaming also commutes
with quantifiers, and the two structures have the same domain, so for
\(Qx\,\alpha_0\) apply the induction hypothesis under every assignment
\(s[x\mapsto a]\) and then the truth clause for \(Q\).
\end{proof}
\end{annexitem}

\begin{annexitem}{annex:p02-w02-p02-037-implicit-agree}{thm:p02-w02-p02-037-implicit-agree}{Implicit definability: two formulations agree}
\begin{proof}
Write \(\bar x=x_1,\ldots,x_n\) and
\(\delta=\forall\bar x\,(P(\bar x)\leftrightarrow P'(\bar x))\).  Because the
\(x_i\) are pairwise distinct, for every structure with domain \(D\), every
assignment \(s\) and every \(\bar a\in D^n\) the assignment
\(s[x_1\mapsto a_1]\cdots[x_n\mapsto a_n]\) sends each \(x_i\) to \(a_i\).
Hence a structure \(\mathcal N\) satisfies \(\delta\) exactly when
\(P^{\mathcal N}=P'^{\mathcal N}\).  For \(n=0\), \(\delta\) is
\(P\leftrightarrow P'\) and the same holds for the two truth values.

\emph{1 implies 2.}  Let \(\mathcal M_1,\mathcal M_2\) be models of
\(T(P)\) with \(\mathcal M_1\mathbin{\upharpoonright}\Lang=\mathcal M_2\mathbin{\upharpoonright}\Lang\).  Let
\(\mathcal N\) be the \(\Lang(P,P')\)-structure with this common
\(\Lang\)-reduct, \(P^{\mathcal N}=P^{\mathcal M_1}\) and
\(P'^{\mathcal N}=P^{\mathcal M_2}\).  Then
\(\mathcal N\mathbin{\upharpoonright}\Lang(P)=\mathcal M_1\), so \(\mathcal N\entails T(P)\) by
\cref{lem:reduct-invariance}.  Also \(\mathcal N^{\mathrm{ren}}=\mathcal M_2\), so
\(\mathcal N\entails\sigma'\) for every \(\sigma\in T(P)\) by
Lemma~\ref{thm:p02-w02-p02-037-renaming}; that is,
\(\mathcal N\entails T(P')\).  By condition~1,
\(\mathcal N\entails\delta\), so
\(P^{\mathcal M_1}=P^{\mathcal N}=P'^{\mathcal N}=P^{\mathcal M_2}\).

\emph{2 implies 1.}  Let \(\mathcal N\entails T(P)\cup T(P')\).  Put
\(\mathcal M_1=\mathcal N\mathbin{\upharpoonright}\Lang(P)\) and
\(\mathcal M_2=\mathcal N^{\mathrm{ren}}\).  Then \(\mathcal M_1\entails T(P)\) by
\cref{lem:reduct-invariance}, and \(\mathcal M_2\entails T(P)\) by
Lemma~\ref{thm:p02-w02-p02-037-renaming}.  Both have the \(\Lang\)-reduct
\(\mathcal N\mathbin{\upharpoonright}\Lang\).  By condition~2,
\(P^{\mathcal N}=P^{\mathcal M_1}=P^{\mathcal M_2}=P'^{\mathcal N}\), so
\(\mathcal N\entails\delta\).  The members of \(T(P)\cup T(P')\) are
sentences, so this is the consequence \(T(P)\cup T(P')\entails\delta\)
under every assignment.
\end{proof}
\end{annexitem}

\begin{annexitem}{annex:p02-w02-p02-037-beth}{thm:p02-w02-p02-037-beth}{Beth definability, renamed-copy form}
\begin{proof}
Write \(\bar x=x_1,\ldots,x_n\).

\paragraph{Explicit implies implicit.}
Suppose \(T(P)\entails\forall\bar x\,(P(\bar x)\leftrightarrow\chi)\) for an
\(\Lang\)-formula \(\chi\).  Let \(\mathcal N\entails T(P)\cup T(P')\).  By
\cref{lem:reduct-invariance}, \(\mathcal N\mathbin{\upharpoonright}\Lang(P)\entails T(P)\),
hence
\(\mathcal N\entails\forall\bar x\,(P(\bar x)\leftrightarrow\chi)\).  By
Lemma~\ref{thm:p02-w02-p02-037-renaming}, \(\mathcal N^{\mathrm{ren}}\entails T(P)\),
hence \(\mathcal N^{\mathrm{ren}}\) satisfies the same sentence; since \(\chi\)
contains no \(P\), its renamed form is
\(\forall\bar x\,(P'(\bar x)\leftrightarrow\chi)\), and the lemma gives
\(\mathcal N\entails\forall\bar x\,(P'(\bar x)\leftrightarrow\chi)\).  Under
every assignment the two biconditionals give
\(P(\bar x)\leftrightarrow P'(\bar x)\), so
\(T(P)\cup T(P')\entails\forall\bar x\,(P(\bar x)\leftrightarrow
P'(\bar x))\).  No choice principle is used.

\paragraph{Implicit implies explicit.}
Suppose \(T(P)\) implicitly defines \(P\) in the sense of
Definition~\ref{def:p02-beth-definability}.  By
Lemma~\ref{thm:p02-w02-p02-037-implicit-agree}, any two models of \(T(P)\)
with the same \(\Lang\)-reduct interpret \(P\) identically.  This is the
hypothesis of Theorem~\ref{thm:p03-beth-definability}, proved in
Annex~\ref{annex:p3.abstract.2}, and of Theorem~\ref{thm:beth}, proved in
Annex~\ref{annex:p7.interp.4}, both under the ambient Choice assumption and
for an arbitrary signature with a relation symbol of arity \(n\geq0\).
Either theorem gives a finite \(T_0\subseteq T(P)\), pairwise distinct
variables \(\bar y=y_1,\ldots,y_n\) and an \(\Lang\)-formula \(\theta\) with
\(\FV(\theta)\subseteq\{y_1,\ldots,y_n\}\) such that \(T_0\) derives
\(\forall\bar y\,(P(\bar y)\leftrightarrow\theta)\).  By soundness,
\cref{thm:fol-soundness},
\[
 T_0\entails\forall\bar y\,(P(\bar y)\leftrightarrow\theta).
 \tag{\(\dagger\)}
\]

If \(n=0\), both quantifier blocks are empty and \(\chi=\theta\) is the
required sentence.  If \(\bar y\) is \(\bar x\), take \(\chi=\theta\).
Otherwise choose pairwise distinct variables \(z_1,\ldots,z_n\) outside
\(\{x_1,\ldots,x_n,y_1,\ldots,y_n\}\cup\operatorname{Var}(\theta)\), and put
\[
 \chi=\exists z_1\cdots\exists z_n
 \Bigl(\bigwedge_{i=1}^{n}z_i=x_i\land
 \exists y_1\cdots\exists y_n
 \bigl(\textstyle\bigwedge_{i=1}^{n}y_i=z_i\land\theta\bigr)\Bigr).
\]
Then \(\chi\) is an \(\Lang\)-formula, since equality is logical.  The inner
existential block binds every free variable of \(\theta\) and leaves only
\(z_1,\ldots,z_n\) free, and the outer block binds these, so
\(\FV(\chi)\subseteq\{x_1,\ldots,x_n\}\).

For a structure \(\mathcal M\) and assignment \(s\), write
\(s(\bar x)=(s(x_1),\ldots,s(x_n))\).  Because the \(z_i\) are distinct from
the \(x_i\) and the \(y_i\), \(\mathcal M,s\entails\chi\) holds iff there
are \(\bar b\) with \(b_i=s(x_i)\) and \(\bar c\) with \(c_i=b_i\) such that
\(\theta\) is satisfied under
\(s[\bar z\mapsto\bar b][\bar y\mapsto\bar c]\).  That assignment agrees
with \(s[\bar y\mapsto s(\bar x)]\) on \(\{y_1,\ldots,y_n\}\supseteq
\FV(\theta)\), so by \cref{lem:fol-coincidence}
\[
 \mathcal M,s\entails\chi
 \quad\Longleftrightarrow\quad
 \mathcal M,s[\bar y\mapsto s(\bar x)]\entails\theta .
\]
Let \(\mathcal M\entails T_0\) and let \(s\) be an assignment.  By
(\(\dagger\)) under the assignment \(s[\bar y\mapsto s(\bar x)]\), the tuple
\(s(\bar x)\) lies in \(P^{\mathcal M}\) iff
\(\mathcal M,s[\bar y\mapsto s(\bar x)]\entails\theta\), that is, iff
\(\mathcal M,s\entails\chi\).  Thus
\(\mathcal M,s\entails P(\bar x)\leftrightarrow\chi\) for every \(s\), and
\[
 T_0\entails\forall\bar x\,(P(\bar x)\leftrightarrow\chi),
 \qquad
 \FV(\chi)\subseteq\{x_1,\ldots,x_n\}.
\]
Since \(T_0\subseteq T(P)\), \(T(P)\) explicitly defines \(P\) over
\(\Lang\) in the sense of Definition~\ref{def:p02-beth-definability}.

\paragraph{Implicit implies explicit: a second proof from the formula form of interpolation.}
This argument uses neither Theorem~\ref{thm:beth} nor
Theorem~\ref{thm:p03-beth-definability}; it uses
Theorem~\ref{thm:p02-w02-p02-037-craig-formula}, compactness and
Lemma~\ref{thm:p02-w02-p02-037-renaming}.  Write
\(\delta=\forall\bar x\,(P(\bar x)\leftrightarrow P'(\bar x))\), a sentence
of \(\Lang(P,P')\), and suppose \(T(P)\cup T(P')\entails\delta\).

\emph{Step 1: a finite part.}  By \cref{thm:fol-compactness}, clause~2, in
the signature \(\Lang(P,P')\), some finite
\(F\subseteq T(P)\cup T(P')\) entails \(\delta\).  Let \(T_0\) consist of
the members of \(F\) that lie in \(T(P)\) together with, for each other
member of \(F\), one \(\sigma\in T(P)\) whose renamed form \(\sigma'\) is
that member.  Only finitely many choices are made, so \(T_0\) is a finite
subset of \(T(P)\), and with \(T_0'=\{\sigma':\sigma\in T_0\}\) we have
\(F\subseteq T_0\cup T_0'\), hence \(T_0\cup T_0'\entails\delta\).  Let
\(\alpha\) be the conjunction of the members of \(T_0\) in a fixed order,
and \(\alpha=\neg\bottom\) if \(T_0\) is empty.  Renaming commutes with the
connectives, so \(\alpha'\) is the conjunction of the members of \(T_0'\)
in the same order.  A structure satisfies \(\alpha\) iff it satisfies
\(T_0\), and \(\alpha'\) iff it satisfies \(T_0'\).

\emph{Step 2: an implication between formulas.}  Put
\[
 \varphi=\alpha\land P(\bar x)\in\operatorname{Fm}(\Lang(P)),
 \qquad
 \psi=\alpha'\to P'(\bar x)\in\operatorname{Fm}(\Lang(P')).
\]
The signatures \(\Lang(P)\) and \(\Lang(P')\) are compatible: their common
symbols are those of \(\Lang\), with the same kinds and arities, because
\(P\) and \(P'\) are distinct and neither belongs to \(\Lang\).  Let
\(\mathcal N\) be an \(\Lang(P,P')\)-structure and \(s\) an assignment
with \(\mathcal N,s\entails\varphi\) and \(\mathcal N,s\entails\alpha'\).
Then \(\mathcal N\) satisfies \(T_0\cup T_0'\), hence \(\delta\), hence
\(P^{\mathcal N}=P'^{\mathcal N}\) by the first paragraph of the proof of
Lemma~\ref{thm:p02-w02-p02-037-implicit-agree}.  Since
\(\mathcal N,s\entails P(\bar x)\), the tuple
\((s(x_1),\ldots,s(x_n))\) lies in \(P'^{\mathcal N}\), that is,
\(\mathcal N,s\entails P'(\bar x)\).  So \(\entails\varphi\to\psi\).

\emph{Step 3: the interpolant.}  Theorem~\ref{thm:p02-w02-p02-037-craig-formula}
gives \(\chi\) with
\[
 \text{(a)}\ \entails\varphi\to\chi,
 \qquad
 \text{(b)}\ \entails\chi\to\psi,
\]
\(\operatorname{voc}(\chi)\subseteq\operatorname{voc}(\varphi)\cap
\operatorname{voc}(\psi)\) and
\(\FV(\chi)\subseteq\FV(\varphi)\cap\FV(\psi)\).  Now
\(\operatorname{voc}(\varphi)\subseteq\Lang\cup\{P\}\) and
\(\operatorname{voc}(\psi)\subseteq\Lang\cup\{P'\}\), so every nonlogical
symbol of \(\chi\) is in \(\Lang\) and \(\chi\) is an \(\Lang\)-formula.
Since \(\alpha\) and \(\alpha'\) are sentences,
\(\FV(\varphi)=\FV(\psi)=\{x_1,\ldots,x_n\}\), so
\(\FV(\chi)\subseteq\{x_1,\ldots,x_n\}\).

\emph{Step 4: renaming back.}  Let \(\mathcal M\) be an
\(\Lang(P)\)-structure and \(s\) an assignment, and let \(\mathcal N\) be
the \(\Lang(P,P')\)-structure with the domain and \(\Lang(P)\)-interpretations
of \(\mathcal M\) and with \(P'^{\mathcal N}=P^{\mathcal M}\).  Then
\(\mathcal N\mathbin{\upharpoonright}\Lang(P)=\mathcal M\) and
\(\mathcal N^{\mathrm{ren}}=\mathcal M\).  By (a) and
\cref{lem:reduct-invariance}, \(\mathcal M,s\entails\varphi\to\chi\).  Put
\(\gamma=\chi\to(\alpha\to P(\bar x))\), an \(\Lang(P)\)-formula.  Since
\(\chi\) contains no \(P\), its renamed form \(\gamma'\) is
\(\chi\to(\alpha'\to P'(\bar x))\), which is \(\chi\to\psi\).  By (b),
\(\mathcal N,s\entails\gamma'\), and
Lemma~\ref{thm:p02-w02-p02-037-renaming} gives
\(\mathcal M,s\entails\gamma\).

\emph{Step 5: the definition.}  Let \(\mathcal M\entails T_0\) and let
\(s\) be an assignment; then \(\mathcal M,s\entails\alpha\).  If
\(\mathcal M,s\entails P(\bar x)\), then \(\mathcal M,s\entails\varphi\), and
Step~4 gives \(\mathcal M,s\entails\chi\).  If
\(\mathcal M,s\entails\chi\), then \(\mathcal M,s\entails\gamma\) gives
\(\mathcal M,s\entails P(\bar x)\).  Hence
\(\mathcal M,s\entails P(\bar x)\leftrightarrow\chi\) for every assignment
\(s\), so \(\mathcal M\entails\forall\bar x\,(P(\bar x)\leftrightarrow\chi)\).
Thus
\[
 T_0\entails\forall\bar x\,(P(\bar x)\leftrightarrow\chi),
 \qquad
 \FV(\chi)\subseteq\{x_1,\ldots,x_n\},
\]
with \(T_0\subseteq T(P)\) finite and \(\chi\) an \(\Lang\)-formula.  For
\(n=0\), \(P(\bar x)\) is the nullary atom, \(\delta\) is
\(P\leftrightarrow P'\), and \(\chi\) is a sentence.  This proof uses the
ambient Choice assumption of the corollary through
\cref{thm:fol-compactness}, which is applied to a possibly infinite theory
in a possibly infinite signature.
\end{proof}
\end{annexitem}
